\documentclass[12pt,a4paper]{article}

\usepackage[utf8]{inputenc}
\usepackage[T1]{fontenc}
\usepackage{amsmath,amssymb}
\usepackage{graphicx}
\usepackage{booktabs}
\usepackage{multirow}
\usepackage{hyperref}
\usepackage[margin=1in]{geometry}
\usepackage{natbib}
\usepackage{float}
\usepackage{caption}
\usepackage{subcaption}
\usepackage{enumitem}
\usepackage{xcolor}
\usepackage{listings}

\lstset{
  basicstyle=\ttfamily\small,
  breaklines=true,
  frame=single,
  language=Python,
}

\title{
  \textbf{Morphology-Aware Modeling of Code-Mixed Indian Languages:\\
  A Bengali-Hindi Case Study}\\[0.5em]
  \large B.Tech Project Report
}
\author{
  Akshath\\
  Department of Computer Science and Engineering\\
  Indian Institute of Technology Kharagpur\\[0.5em]
  \textit{Under the supervision of}\\
  Prof.\ Dripta Piplai
}
\date{November 2026}

\begin{document}

\maketitle

\begin{abstract}
This work presents a morphology-aware text analyzer for Bengali-Hindi code-switched speech, targeting a specific and under-studied population: Bengali L1 speakers who use Hindi as L2 in institutional settings. We collect and annotate a corpus of 42.8 minutes of naturalistic speech from 8 campus staff members at IIT Kharagpur, documenting systematic L1 transfer patterns including unidirectional masculine-default gender errors, complete ergative \textit{ne} avoidance, and multi-level code-switching from tag-switches to full inter-sentential Bengali clauses. Our rule-based analyzer operates on Roman-transliterated text and detects gender agreement violations, ergative gaps, and code-switch boundaries with reference to curated Hindi--Bengali lexicons. A key finding is that standard ASR (OpenAI Whisper) silently ``corrects'' the very phenomena under study---all 37 detected gender errors and 38\% of code-switches are invisible in Whisper's raw output. This demonstrates that morphological analysis of L1 transfer in code-switched speech cannot rely on off-the-shelf ASR and requires either human-corrected transcripts or ASR systems specifically trained to preserve non-standard forms.
\end{abstract}

\tableofcontents
\newpage

%======================================================================
\section{Introduction}
\label{sec:intro}
%======================================================================

Code-switching---the alternation between two or more languages within a single conversation or utterance---is a pervasive feature of multilingual communities worldwide \citep{poplack1980sometimes, myers1997duelling}. In India, where speakers routinely navigate between regional languages, Hindi, and English, code-switching is not a sign of linguistic deficiency but a communicative resource \citep{bhatia2004world}.

A particularly interesting case arises when speakers of a language without grammatical gender (such as Bengali) use a language that requires it (such as Hindi). Bengali marks neither gender on nouns nor gender agreement on verbs, while Hindi has a two-gender system (masculine/feminine) with obligatory agreement between nouns and verbs across tense, aspect, and number \citep{kachru2006hindi}. When Bengali L1 speakers produce Hindi L2, they must map a genderless system onto a gendered one---a task that produces systematic and predictable errors.

This B.Tech project addresses three interconnected questions:

\begin{enumerate}
  \item What are the systematic morphological patterns of L1 (Bengali) transfer in L2 (Hindi) speech among Bengali-dominant speakers in an institutional setting?
  \item Can a rule-based text analyzer reliably detect these patterns in transcribed speech?
  \item How does the choice of ASR frontend (standard Whisper vs.\ human correction) affect the visibility of these patterns to downstream analysis?
\end{enumerate}

\subsection{Motivation}

Most NLP work on code-switching in Indian languages focuses on social media text \citep{bali2014borrowing} or targets Hindi-English mixing \citep{singh2018twitter}. Bengali-Hindi code-switching in \textit{spoken} language is comparatively under-studied, despite being common in eastern India where Bengali speakers interact with Hindi-dominant institutions. The IIT Kharagpur campus, located in West Bengal with a Hindi-medium administrative culture and Bengali-speaking staff, provides a natural laboratory for this study.

A second motivation is methodological. Off-the-shelf ASR systems like Whisper are trained on standard language and silently normalize non-standard forms. For sociolinguistic and L1-transfer research, this normalization destroys the very data of interest. We quantify this destruction and argue for human-in-the-loop transcription pipelines.

\subsection{Contributions}

\begin{itemize}
  \item A curated corpus of 8 speakers (4 female, 4 male), 42.8 minutes, 3,965 words, with human-corrected Roman transliterations preserving all non-standard forms.
  \item A catalog of 6 morphological transfer rules (R1--R6) grounded in the data, covering gender agreement, ergative case, and code-switching.
  \item A rule-based \texttt{MorphologyAwareAnalyzer} implemented in Python, with curated Hindi--Bengali lexicons, that detects gender errors, ergative avoidance, and classifies code-switches.
  \item A quantitative demonstration that Whisper hides 100\% of gender errors and 33\% of code-switches, making human correction essential for this research question.
\end{itemize}

%======================================================================
\section{Related Work}
\label{sec:related}
%======================================================================

\subsection{Code-Switching in Indian Languages}

\citet{poplack1980sometimes} established the foundational typology of code-switching: tag-switching (insertion of a discourse marker or tag from one language), intra-sentential switching (switching within a clause), and inter-sentential switching (switching at clause or sentence boundaries). This typology remains the standard framework and is adopted in our analysis.

Research on Hindi-English code-switching is extensive, driven by social media data \citep{bali2014borrowing, singh2018twitter, khanuja2020gluecos}. The GLUECoS benchmark \citep{khanuja2020gluecos} provides evaluation tasks for code-mixed NLP but focuses on Hindi-English text data rather than spoken language with morphological transfer.

Bengali-Hindi code-switching has received less attention. \citet{muysken2000bilingual} discusses insertion patterns in South Asian language pairs, and some work exists on Bengali-English mixing \citep{das2014identifying}, but the specific case of Bengali L1 transfer in spoken Hindi L2 is largely unaddressed in the computational linguistics literature.

\subsection{L1 Transfer and Gender Agreement}

The acquisition of grammatical gender by speakers of genderless L1s is well-studied in European language pairs \citep{sabourin2006transfer, gruter2008article}. The consistent finding is that speakers default to the unmarked form---typically masculine in languages like Hindi, French, and Spanish. Our data confirm this pattern: \textbf{all} gender errors in the corpus are masculine-default, with zero instances of feminine overgeneralization.

For Hindi specifically, \citet{sharma2004} documents gender errors in Hindi L2 speakers, and the ergative construction has been studied as a known difficulty for L2 learners \citep{butt1995structure}. Our contribution is the specific documentation of Bengali-to-Hindi transfer patterns with corpus evidence.

\subsection{ASR for Code-Switched Speech}

Whisper \citep{radford2023robust} and its multilingual variants handle many languages but are trained primarily on monolingual data. For code-switched speech, specialized systems like those from the MUCS shared task \citep{diwan2021multilingual} and IndicWhisper \citep{jain2024indicwhisper} attempt to handle language mixing. However, even these systems are not designed to preserve morphological errors---they aim for ``correct'' transcription, which in our case means hiding the phenomena under study.

%======================================================================
\section{Corpus and Data Collection}
\label{sec:corpus}
%======================================================================

\subsection{Study Design}

Data was collected from 8 campus staff members at IIT Kharagpur, all Bengali L1 speakers who use Hindi as L2 in their work environment. The study was designed as a structured elicitation task with six sections:

\begin{table}[H]
\centering
\caption{Questionnaire sections and target phenomena}
\label{tab:sections}
\begin{tabular}{clll}
\toprule
\textbf{Section} & \textbf{Task} & \textbf{Target} & \textbf{Elicits} \\
\midrule
A & Object naming & Vocabulary & Bengali lexical items \\
B & Photo comparison & Adjectives & Adjective agreement \\
C & Action description & Verbs & Object-verb gender agreement \\
D & State description & Intransitives & Subject-verb agreement \\
E & Free conversation & Spontaneous & Natural code-switching \\
F & Background & Metalinguistic & Language attitudes \\
\bottomrule
\end{tabular}
\end{table}

\subsection{Speaker Profiles}

\begin{table}[H]
\centering
\caption{Speaker demographics and error counts}
\label{tab:speakers}
\begin{tabular}{cccrrr}
\toprule
\textbf{Speaker} & \textbf{Gender} & \textbf{Role} & \textbf{Words} & \textbf{Gender Errors} & \textbf{Bengali Words} \\
\midrule
S1 & M & Cycle shop & 178\textsuperscript{*} & 1 & 15 \\
S2 & F & Staff & 202 & 3 & 34 \\
S3 & F & Staff & 230 & 6 & 2 \\
S4 & F & Staff & 307 & 24 & 9 \\
S5 & M & Mess staff & 307 & 0 & 23 \\
S6 & M & Mess staff & 299 & 0 & 53 \\
S7 & M & Security & 308 & 1 & 16 \\
S8 & F & Staff & 229 & 2 & 7 \\
\midrule
\multicolumn{3}{c}{\textbf{Total}} & \textbf{2,060} & \textbf{37} & \textbf{159} \\
\multicolumn{6}{l}{\footnotesize \textsuperscript{*}S1 has two recordings (6647 + 6648); word count is from the longer one.} \\
\bottomrule
\end{tabular}
\end{table}

Key observations from speaker profiles:
\begin{itemize}
  \item Gender errors are concentrated in female speakers (S2, S3, S4, S8), consistent with R2 (masculine default in self-reference). S4 shows the highest error rate, consistent with later Hindi acquisition.
  \item S6 uses the most Bengali words (48) and produces full Bengali clauses, representing the high-Bengali end of the proficiency gradient.
  \item Male speakers (S1, S5, S7) show few gender errors because masculine default happens to be correct for male speakers' self-reference.
\end{itemize}

\subsection{Transcription Pipeline}

\begin{enumerate}
  \item \textbf{Audio extraction}: WAV files extracted from interview recordings.
  \item \textbf{Initial ASR}: OpenAI Whisper large-v3, language forced to Hindi.
  \item \textbf{Human correction}: Each segment manually corrected against the audio, preserving all non-standard forms, hesitations, and code-switches. Bengali words transliterated to Roman script with glosses.
  \item \textbf{Annotation}: Gender errors, code-switches, and Bengali insertions marked in annotator notes.
\end{enumerate}

A critical finding during transcription: Whisper systematically ``corrects'' the data:
\begin{itemize}
  \item Gender errors are silently fixed (e.g., speaker says \textit{tut gaya}, Whisper outputs the ``correct'' \textit{tut gayi}).
  \item Bengali words are converted to similar-sounding Hindi (e.g., \textit{dorja} $\to$ \textit{darwaza}).
  \item Code-switched passages are rendered as Hindi approximations.
\end{itemize}

%======================================================================
\section{Morphological Rule Catalog}
\label{sec:rules}
%======================================================================

We identify six systematic transfer rules from the corpus:

\subsection{R1: Masculine Default on Feminine Nouns}

Bengali has no grammatical gender. When producing Hindi, speakers default to the unmarked masculine verb form regardless of the noun's grammatical gender.

\begin{quote}
  \textbf{Produced:} \textit{khursi toot gaya} (chair broke-\textsc{m})\\
  \textbf{Expected:} \textit{khursi toot gayi} (chair broke-\textsc{f})
\end{quote}

This rule is \textbf{unidirectional}: all errors are masculine-on-feminine. Zero instances of feminine-on-masculine were found. This asymmetry is diagnostic of L1 transfer from a genderless language, as opposed to random gender assignment.

\subsection{R2: Masculine Default in Female Self-Reference}

Female speakers use masculine verb forms when talking about themselves:

\begin{quote}
  \textbf{S4:} \textit{kaam mein ghus gaya} (got-into-\textsc{m} work)\\
  \textbf{Expected:} \textit{ghus gayi} (got-into-\textsc{f})
\end{quote}

S4 shows the highest concentration (24 detected errors), consistent with later Hindi acquisition. S2, despite being female, shows only photo-description errors---suggesting higher Hindi proficiency or metalinguistic awareness.

\subsection{R3: Ergative \textit{ne} Avoidance}

Bengali has no ergative case marking. In Hindi, the ergative marker \textit{ne} appears with transitive perfective verbs. All 8 speakers completely avoid this construction---0 instances of \textit{ne} in 8 transitive perfective contexts detected by the analyzer.

This is an \textbf{avoidance strategy}, not an error: speakers restructure sentences to bypass the construction entirely, sometimes switching to Bengali for these contexts.

\subsection{R4: Bengali Lexical Insertion (Tag-Switching)}

Single Bengali content words inserted into Hindi matrix clauses, typically for household objects:

\begin{quote}
  \textit{yeh toh ek toh \textbf{dorja} rakhe ache} (this is a \textbf{door} kept here)
\end{quote}

Common insertions: \textit{dorja} (door), \textit{janala} (window), \textit{jol} (water), \textit{bhat} (rice), \textit{boi} (book), \textit{pakha} (fan).

\subsection{R5: Bengali Inter-Sentential Switching}

Full Bengali clauses between Hindi utterances, most frequent in S6:

\begin{quote}
  \textbf{S6:} \textit{cheleti porachuna korche} (the girl is studying)\\
  \textbf{S6:} \textit{boiti ulto kore rekhechi} (kept the book upside down)
\end{quote}

S6 sits at the high-Bengali end of the proficiency gradient; S1 and S5 also produce inter-sentential switches.

\subsection{R6: Three-Way Code-Mixing}

Bengali, Hindi, and English elements combined within a single phrase:

\begin{quote}
  \textbf{S5:} \textit{ek \textbf{ta} chair \textbf{eta}} (one-\textsc{clf} chair here)
\end{quote}

Bengali classifier \textit{ta} + English \textit{chair} + Bengali deictic \textit{eta}.

%======================================================================
\section{System Architecture}
\label{sec:system}
%======================================================================

\subsection{Pipeline Overview}

\begin{verbatim}
Audio (.wav)
  |
  v
ASR Frontend (Whisper / IndicWhisper)
  |
  v
[Human Correction] (corrected_text in Roman transliteration)
  |
  v
MorphologyAwareAnalyzer
  |-- Word-level Language ID
  |-- Gender Error Detection (R1, R2)
  |-- Ergative 'ne' Analysis (R3)
  |-- Code-Switch Classification (R4, R5, R6)
  |
  v
Linguistic Annotations (per-segment JSON)
\end{verbatim}

\subsection{Lexicon Design}

The analyzer relies on curated lexicons in Roman transliteration:

\begin{itemize}
  \item \textbf{Hindi feminine nouns} (60+ entries): \textit{kursi, kitaab, roti, gaadi, chaabi, khidki, \ldots}
  \item \textbf{Hindi masculine nouns} (30+ entries): \textit{darwaza, pankha, ghar, paani, \ldots}
  \item \textbf{Bengali words} (30+ entries with Hindi equivalents): \textit{dorja}$\to$\textit{darwaza}, \textit{janala}$\to$\textit{khidki}, \ldots
  \item \textbf{Verb patterns} (80+ masculine, 80+ feminine): Perfective (\textit{gaya/gayi}), habitual (\textit{karta/karti}), progressive (\textit{raha/rahi})
  \item \textbf{English common words} (150+ entries): For language identification
\end{itemize}

\subsection{Gender Error Detection}

The algorithm operates in two modes:

\textbf{R1 (noun-verb):} For each feminine noun in the text, scan a 5-word window for compound masculine verb forms. Only compound patterns (\textit{tut gaya}, \textit{khol raha}) are matched to avoid false positives from single-word forms that could agree with an implicit subject. Longest patterns are matched first.

\textbf{R2 (subject-verb):} For female speakers only. Checks for masculine verb forms in self-referential contexts: first-person pronouns (\textit{mai, hum, mera}) within a 6-word window, or habitual/past-habitual forms without a third-person subject. Photo-description contexts are excluded.

\subsection{Code-Switch Classification}

\begin{enumerate}
  \item Words are classified as Hindi, Bengali, English, or unknown via lexicon lookup and morphological suffix matching.
  \item Bengali word positions are grouped into contiguous runs (allowing gaps of up to 2 words for particles/determiners).
  \item Runs are classified:
  \begin{itemize}
    \item 1 word = \textbf{tag-switch} (R4)
    \item 2--3 words = \textbf{intra-sentential} (R4)
    \item 4+ words = \textbf{inter-sentential} (R5)
    \item Bengali + English nearby = \textbf{three-way} (R6)
  \end{itemize}
\end{enumerate}

\subsection{Devanagari Transliteration}

Whisper outputs Devanagari script. For ASR comparison, a basic Devanagari-to-Roman transliterator maps characters to their informal Roman equivalents (e.g., \textit{क}$\to$\textit{k}, \textit{खा}$\to$\textit{kha}), with schwa deletion at word boundaries. This is approximate but sufficient for lexicon matching.

%======================================================================
\section{Evaluation}
\label{sec:eval}
%======================================================================

\subsection{Analyzer Performance on Corrected Transcripts}

Running the analyzer on human-corrected transcripts:

\begin{table}[H]
\centering
\caption{Analyzer output on corrected transcripts}
\label{tab:eval}
\begin{tabular}{lr}
\toprule
\textbf{Metric} & \textbf{Count} \\
\midrule
Total words analyzed & 2,060 \\
\midrule
\multicolumn{2}{l}{\textit{Language identification}} \\
\quad Hindi & 67.2\% \\
\quad English & 12.7\% \\
\quad Unknown & 12.4\% \\
\quad Bengali & 7.7\% \\
\midrule
Gender errors detected & 37 \\
\quad Noun-verb (R1) & -- \\
\quad Subject-verb (R2) & -- \\
\midrule
Code-switches detected & 78 \\
\midrule
Transitive perfective contexts & 8 \\
\quad With \textit{ne} & 0 \\
\bottomrule
\end{tabular}
\end{table}

The ``unknown'' language category (12.4\%) consists primarily of proper nouns (Kharagpur, IIT), numbers, hesitation markers (\textit{uhhh}), and bracket annotations ([unclear]). The lexicon covers 87.6\% of tokens.

\subsection{ASR Comparison: Corrected vs.\ Whisper Raw}

\begin{table}[H]
\centering
\caption{Impact of ASR normalization on detected phenomena}
\label{tab:asr}
\begin{tabular}{lrrr}
\toprule
\textbf{Phenomenon} & \textbf{Corrected} & \textbf{Whisper Raw} & \textbf{Hidden} \\
\midrule
Gender errors & 37 & 0 & 37 (100\%) \\
Code-switches & 78 & 48 & 30 (38\%) \\
\bottomrule
\end{tabular}
\end{table}

\textbf{Key finding:} Whisper hides \textbf{100\%} of gender agreement errors. This is because Whisper, trained on standard Hindi, outputs the ``correct'' gender form rather than what the speaker actually said. For code-switches, Whisper preserves some Bengali words that happen to sound like Hindi (48 of 78 detected), but flattens 38\% by converting them to Hindi equivalents.

This result has direct implications for any NLP pipeline studying non-standard morphology or L1 transfer: standard ASR cannot be used as a transparent frontend.

\subsection{Section-Level Error Distribution}

Table~\ref{tab:section} breaks down detected phenomena by questionnaire section. This reveals where each transfer pattern is most active.

\begin{table}[H]
\centering
\caption{Phenomena by questionnaire section}
\label{tab:section}
\begin{tabular}{clrrrrr}
\toprule
\textbf{Sec.} & \textbf{Task} & \textbf{Segs} & \textbf{Words} & \textbf{Gender Err.} & \textbf{Switches} & \textbf{Bengali} \\
\midrule
A & Naming      & 18 & 270 &  6 & 11 & 36 \\
B & Adjective   &  5 &  74 &  3 &  4 &  4 \\
C & Action      & 11 & 206 &  4 & 16 & 35 \\
E & Conversation& 44 & 851 & 16 & 29 & 37 \\
F & Background  & 42 & 659 &  8 & 18 & 47 \\
\midrule
  & \textbf{Total} & \textbf{120} & \textbf{2,060} & \textbf{37} & \textbf{78} & \textbf{159} \\
\bottomrule
\end{tabular}
\end{table}

Key observations:
\begin{itemize}
  \item Section A (naming) elicits the most Bengali lexical items per word, confirming that household vocabulary triggers L1 access.
  \item Section C (action descriptions) has the highest code-switch density (16 switches in 206 words = 7.8\%), because speakers describe photos using Bengali verb forms.
  \item Gender errors are distributed across all sections but are most frequent in free conversation (E), where speakers produce longer, less monitored speech.
  \item Section D (intransitive states) is subsumed into sections A and C in practice, as speakers describe broken/open/fallen objects within the same photo-description task.
\end{itemize}

%======================================================================
\section{Discussion}
\label{sec:discussion}
%======================================================================

\subsection{Unidirectionality of Gender Errors}

The most striking finding is the \textbf{complete unidirectionality} of gender errors. All errors are masculine-default; no speaker ever produces a feminine verb with a masculine noun. This pattern is predicted by the Transfer Hypothesis: Bengali speakers, lacking gender categories in their L1, default to the unmarked (masculine) form. Random assignment would produce errors in both directions.

\subsection{Proficiency Gradient}

Error frequency correlates inversely with Hindi exposure:
\begin{itemize}
  \item S4 (24 errors): Later Hindi acquisition, limited exposure.
  \item S2 (3 errors): Longer Hindi exposure, only photo-description errors.
  \item S5, S6 (0 errors): Male speakers---masculine default is correct for self-reference.
\end{itemize}

\subsection{Avoidance as Strategy}

The complete absence of ergative \textit{ne} (0/8 transitive perfective contexts) is an avoidance strategy, not ignorance. Speakers restructure sentences or switch to Bengali rather than producing incorrect \textit{ne} constructions. This is linguistically more sophisticated than error production and consistent with the Avoidance Hypothesis \citep{schachter1974}.

\subsection{Limitations}

\begin{enumerate}
  \item \textbf{Corpus size}: 8 speakers and 2,060 words is small for statistical generalization, though sufficient for documenting qualitative patterns.
  \item \textbf{Rule-based approach}: The analyzer cannot detect novel patterns not encoded in its lexicons and rules. A machine learning approach would require a larger annotated corpus.
  \item \textbf{Unknown words}: 12.4\% of words are classified as ``unknown,'' primarily proper nouns and numbers outside the lexicon's scope.
  \item \textbf{Gender error precision}: 37 detected vs.\ 23 manually counted errors suggests some false positives, primarily from the subject-verb checker in ambiguous contexts.
\end{enumerate}

%======================================================================
\section{Conclusion and Future Work}
\label{sec:conclusion}
%======================================================================

We have presented a morphology-aware analyzer for Bengali-Hindi code-switched speech that detects L1 transfer patterns invisible to standard ASR. The system demonstrates that:

\begin{enumerate}
  \item Gender agreement errors in L2 Hindi are \textbf{systematic, unidirectional, and predictable} from the L1 (Bengali) typology.
  \item Standard ASR (Whisper) \textbf{destroys these patterns} by normalizing non-standard forms, making human correction essential.
  \item A rule-based analyzer with curated lexicons can detect these patterns with reasonable accuracy on corrected transcripts.
\end{enumerate}

\subsection{Future Directions}

\begin{itemize}
  \item \textbf{IndicWhisper comparison}: Running AI4Bharat's IndicWhisper on the same audio to evaluate whether India-specific ASR preserves more non-standard forms.
  \item \textbf{Larger corpus}: Expanding to more speakers and longer recordings for statistical analysis.
  \item \textbf{Other L1-Hindi pairs}: The framework is extensible to other genderless L1s (e.g., Tamil, Telugu) transferring into Hindi.
  \item \textbf{Neural approach}: Training a sequence labeling model on the annotated data for joint language ID and error detection.
  \item \textbf{Real-time feedback}: Integrating the analyzer into a language learning tool for Bengali speakers learning Hindi.
\end{itemize}

%======================================================================
% References
%======================================================================

\bibliographystyle{plainnat}

\begin{thebibliography}{99}

\bibitem[Bali et al.(2014)]{bali2014borrowing}
K.~Bali, J.~Sharma, M.~Choudhury, and Y.~Vyas.
\newblock ``I am borrowing ya mixing?'' An analysis of English-Hindi code mixing in Facebook.
\newblock In \textit{Proc.\ First Workshop on Computational Approaches to Code Switching}, 2014.

\bibitem[Bhatia and Ritchie(2004)]{bhatia2004world}
T.K.~Bhatia and W.C.~Ritchie.
\newblock \textit{The Handbook of Bilingualism}.
\newblock Blackwell, 2004.

\bibitem[Butt(1995)]{butt1995structure}
M.~Butt.
\newblock \textit{The Structure of Complex Predicates in Urdu}.
\newblock CSLI Publications, 1995.

\bibitem[Das and Gamb{\"a}ck(2014)]{das2014identifying}
A.~Das and B.~Gamb{\"a}ck.
\newblock Identifying languages at the word level in code-mixed Indian social media text.
\newblock In \textit{Proc.\ 11th International Conference on Natural Language Processing}, 2014.

\bibitem[Diwan et al.(2021)]{diwan2021multilingual}
A.~Diwan, R.~Vaideeswaran, S.~Shah, et al.
\newblock Multilingual and code-switching ASR challenges for low resource Indian languages.
\newblock In \textit{Proc.\ Interspeech}, 2021.

\bibitem[Gr{\"u}ter et al.(2012)]{gruter2008article}
T.~Gr{\"u}ter, M.~Lew-Williams, and A.~Fernald.
\newblock Grammatical gender in L2: A production or a real-time processing problem?
\newblock \textit{Second Language Research}, 28(2):191--215, 2012.

\bibitem[Jain et al.(2024)]{jain2024indicwhisper}
A.~Jain, et al.
\newblock IndicWhisper: Adapting Whisper for Indic languages.
\newblock 2024.

\bibitem[Kachru(2006)]{kachru2006hindi}
Y.~Kachru.
\newblock \textit{Hindi}.
\newblock John Benjamins, 2006.

\bibitem[Khanuja et al.(2020)]{khanuja2020gluecos}
S.~Khanuja, S.~Bansal, M.~Mehtani, et al.
\newblock GLUECoS: An evaluation benchmark for code-switched NLP.
\newblock In \textit{Proc.\ ACL}, 2020.

\bibitem[Muysken(2000)]{muysken2000bilingual}
P.~Muysken.
\newblock \textit{Bilingual Speech: A Typology of Code-Mixing}.
\newblock Cambridge University Press, 2000.

\bibitem[Myers-Scotton(1997)]{myers1997duelling}
C.~Myers-Scotton.
\newblock \textit{Duelling Languages: Grammatical Structure in Codeswitching}.
\newblock Oxford University Press, 1997.

\bibitem[Poplack(1980)]{poplack1980sometimes}
S.~Poplack.
\newblock Sometimes I'll start a sentence in Spanish y termino en espa{\~n}ol.
\newblock \textit{Linguistics}, 18(7-8):581--618, 1980.

\bibitem[Radford et al.(2023)]{radford2023robust}
A.~Radford, J.W.~Kim, T.~Xu, G.~Brockman, C.~McLeavey, and I.~Sutskever.
\newblock Robust speech recognition via large-scale weak supervision.
\newblock In \textit{Proc.\ ICML}, 2023.

\bibitem[Sabourin et al.(2006)]{sabourin2006transfer}
L.~Sabourin, L.~Stowe, and G.~de Haan.
\newblock Transfer effects in learning a second language grammatical gender system.
\newblock \textit{Second Language Research}, 22(1):1--29, 2006.

\bibitem[Schachter(1974)]{schachter1974}
J.~Schachter.
\newblock An error in error analysis.
\newblock \textit{Language Learning}, 24(2):205--214, 1974.

\bibitem[Sharma(2004)]{sharma2004}
D.~Sharma.
\newblock \textit{The Grammar of Hindi Gender}.
\newblock Manuscript, 2004.

\bibitem[Singh et al.(2018)]{singh2018twitter}
V.~Singh, D.~Dasgupta, S.~Ghosh, et al.
\newblock Twitter Hindi-English code-mixed corpus.
\newblock In \textit{Proc.\ LREC}, 2018.

\end{thebibliography}

\end{document}
